\subsection{発展的な読み物}

SWEBOK 03章で紹介される発展的な読み物を、日本語で読む目的に合わせて整理する。

\par\smallskip\noindent\refstepcounter{table}\label{tab:ch03-09}\textbf{表 \thetable: 第03章 ソフトウェア設計 表09}\par\smallskip
\begin{tabularx}{\linewidth}{L{52mm}Y}
\toprule
\textbf{文献} & \textbf{読むと得られること} \\
\midrule
Brooks, \textit{The Design of Design} & ソフトウェア工学の先駆者による設計論である。設計活動、設計判断、事例を広く学べる。 \\
Budgen, \textit{Software Design} & ソフトウェア設計を体系的に学ぶための中心的な教科書である。構造、振る舞い、設計方法、評価を幅広く扱う。 \\
Gamma et al., \textit{Design Patterns} & 再利用可能なオブジェクト指向設計パターンの古典である。設計語彙を学ぶのに有効である。 \\
Booch, Rumbaugh and Jacobson, \textit{The Unified Modeling Language User Guide} & UMLを使った構造・振る舞いの表現を学べる。設計記述の記法を理解するのに役立つ。 \\
Sommerville, \textit{Software Engineering} & 設計だけでなく、要求、実装、テスト、品質、分散システムなどを含むソフトウェア工学全体を学べる。 \\
Parnas, “On the Criteria To Be Used In Decomposing Systems Into Modules” & モジュール分解と情報隠蔽の古典的論文である。設計原則の背景を学べる。 \\
IEEE Std 7000-2021 & 倫理的関心をシステム設計で扱うための標準である。AIや自律システムの設計で参考になる。 \\
\bottomrule
\end{tabularx}
